%%======================================================================
%%  PPN-1 v2: Literature Extraction & Convention Fixing
%%            for SCT Post-Newtonian Solar System Tests
%%
%%  Comprehensive extraction of all equations, conventions, and
%%  benchmarks from the primary literature needed for the PPN-1
%%  derivation.  This document feeds into PPN1_derivation.tex.
%%
%%  RESEARCH DOCUMENT: formulas extracted from cited papers.
%%  No original derivations --- gaps noted explicitly.
%%======================================================================
\documentclass[11pt,a4paper]{article}

%% ---- packages -------------------------------------------------------
\usepackage[utf8]{inputenc}
\usepackage[T1]{fontenc}
\usepackage{lmodern}
\usepackage{amsmath,amssymb,amsthm,mathtools}
\usepackage[margin=2.5cm]{geometry}
\usepackage{booktabs}
\usepackage{enumitem}
\usepackage{hyperref}
\usepackage{xcolor}
\usepackage{fancyhdr}
\usepackage{longtable}
\usepackage{cite}

%% ---- colors ---------------------------------------------------------
\definecolor{sctblue}{RGB}{0,51,153}
\definecolor{verified}{RGB}{0,128,0}

%% ---- theorem-like environments --------------------------------------
\theoremstyle{plain}
\newtheorem{theorem}{Theorem}[section]
\newtheorem{proposition}[theorem]{Proposition}
\newtheorem{lemma}[theorem]{Lemma}
\newtheorem{corollary}[theorem]{Corollary}
\theoremstyle{definition}
\newtheorem{definition}[theorem]{Definition}
\newtheorem{convention}[theorem]{Convention}
\newtheorem{benchmark}[theorem]{Benchmark}
\newtheorem{finding}[theorem]{Finding}
\theoremstyle{remark}
\newtheorem{remark}[theorem]{Remark}

%% ---- custom commands ------------------------------------------------
\newcommand{\Tr}{\operatorname{Tr}}
\newcommand{\tr}{\operatorname{tr}}
\newcommand{\Id}{\mathbb{1}}
\newcommand{\Ricsq}{R_{\mu\nu}R^{\mu\nu}}
\newcommand{\Weylsq}{C_{\mu\nu\rho\sigma}C^{\mu\nu\rho\sigma}}
\newcommand{\dd}{\mathrm{d}}
\newcommand{\ii}{\mathrm{i}}
\newcommand{\erfi}{\operatorname{erfi}}
\DeclareMathOperator{\sgn}{sgn}
\DeclareMathOperator{\erf}{erf}

%% ---- annotation commands --------------------------------------------
\newcommand{\fromlit}[1]{\textcolor{sctblue}{\textsf{[#1]}}}
\newcommand{\oursign}{\textcolor{teal!80!black}{\textsf{[our convention]}}}
\newcommand{\gap}{\textcolor{red!80!black}{\textsf{[GAP]}}}
\newcommand{\needscheck}{\textcolor{orange!80!black}{\textsf{[needs check]}}}

%% ---- verification annotations ---------------------------------------
\newcommand{\verified}[1]{%
  \par\smallskip\noindent
  \colorbox{green!10}{\parbox{\dimexpr\linewidth-2\fboxsep}{%
    \textcolor{verified}{\textbf{[Verified]}} #1}}%
  \par\smallskip}

\newcommand{\signcorrection}[1]{%
  \par\smallskip\noindent
  \colorbox{red!10}{\parbox{\dimexpr\linewidth-2\fboxsep}{%
    \textcolor{red!60!black}{\textbf{[Important note]}} #1}}%
  \par\smallskip}

\newcommand{\crosscheck}[1]{%
  \par\smallskip\noindent
  \colorbox{blue!10}{\parbox{\dimexpr\linewidth-2\fboxsep}{%
    \textcolor{sctblue}{\textbf{[Cross-check]}} #1}}%
  \par\smallskip}

%% ---- page style -----------------------------------------------------
\pagestyle{fancy}
\fancyhf{}
\fancyhead[L]{\small\textsc{SCT Theory --- Literature Extraction}}
\fancyhead[R]{\small\textsc{PPN-1 v2}}
\fancyfoot[C]{\thepage}

%% ---- hyperref -------------------------------------------------------
\hypersetup{
  colorlinks=true,
  linkcolor=sctblue,
  citecolor=verified,
  urlcolor=sctblue,
  pdftitle={PPN-1 v2: Literature Extraction for SCT Post-Newtonian Tests},
  pdfauthor={David Alfyorov},
}

%% ---- title ----------------------------------------------------------
\title{\textbf{PPN-1 v2: Literature Extraction \& Convention Fixing\\[4pt]
       for SCT Post-Newtonian Solar System Tests}\\[8pt]
       \large Comprehensive Extraction of the 10-Parameter PPN Formalism,\\
       Higher-Derivative Gravity Benchmarks, Nonlocal Potentials,\\
       and Experimental Constraints}
\author{David Alfyorov}
\date{April 2026 --- Working Document}

%%======================================================================
\begin{document}
\maketitle

\begin{abstract}
This document provides a complete, self-contained extraction of all
equations, sign conventions, and experimental benchmarks from the
primary literature needed to derive the PPN parameters of the SCT
spectral action and compare them against solar-system and laboratory
constraints.  Seven sub-tasks are addressed:
(L-1)~the full 10-parameter PPN formalism of Will~(2014);
(L-2)~Stelle's higher-derivative gravity as a local benchmark;
(L-3)~nonlocal gravity potentials from Edholm--Koshelev--Mazumdar~(2016)
and Biswas--Mazumdar--Siegel~(2006);
(L-4)~the Brans--Dicke scalar-tensor limit as a structural cross-check;
(L-5)~the Barnes--Rivers projector algebra for static sources;
(L-6)~a compilation of experimental bounds from Cassini, MESSENGER,
E\"ot-Wash, LLR, and MICROSCOPE;
(L-7)~the van~Dam--Veltman--Zakharov discontinuity and why SCT avoids it.
Every formula is cited to its primary source with equation number.
No original derivations are performed; all new results are deferred
to \texttt{PPN1\_derivation.tex}.

\medskip\noindent
\textbf{Purpose:} Literature extraction feeding into the PPN-1 v2
derivation pipeline.
\end{abstract}

\tableofcontents
\newpage

%%======================================================================
\section{Master Convention Block}\label{sec:master-conventions}
%%======================================================================

All conventions in this document are fixed to match the existing SCT
derivation chain (NT-1 through NT-4a).  Any conflict with the
literature source is noted explicitly.

\begin{convention}[Metric signature]\label{conv:signature}
\oursign\;
We work in the $(-,+,+,+)$ Lorentzian signature for physical
predictions.  The SCT spectral action is formulated in Euclidean
signature with continuation to Lorentzian at the end.  The static
weak-field metric is:
\begin{equation}\label{eq:static-metric}
  \dd s^2 = -(1 + 2\Phi(r))\,\dd t^2
  + (1 - 2\Psi(r))\,\delta_{ij}\,\dd x^i\dd x^j,
\end{equation}
where $\Phi < 0$ and $\Psi < 0$ for an attractive source.  In GR:
$\Phi = \Psi = -GM/r$.
\end{convention}

\begin{convention}[Unit system]\label{conv:units}
We use natural units $\hbar = c = 1$ for all theoretical formulas.
Physical predictions restore SI units via $\hbar c = 197.3$~MeV$\cdot$fm,
$G = 6.674 \times 10^{-11}$~m$^3$kg$^{-1}$s$^{-2}$,
$M_{\mathrm{Pl}} = \sqrt{\hbar c / (8\pi G)} = 2.435 \times 10^{18}$~GeV
(reduced Planck mass).
\end{convention}

\begin{convention}[PPN sign mapping]\label{conv:ppn-sign}
Will~(2014) uses $U = GM/r > 0$ for the Newtonian potential.
Our convention uses $\Phi = -U < 0$.  The mapping is:
\begin{equation}\label{eq:sign-map}
  \Phi = -U, \qquad
  \Psi = -\gamma_{\mathrm{PPN}}\, U \quad\text{(at Newtonian order)}.
\end{equation}
Therefore $\gamma_{\mathrm{PPN}} = \Psi/\Phi$ with the signs cancelling.
\end{convention}

\signcorrection{The sign convention $\Phi < 0$ for attraction
differs from the PPN convention $U > 0$.  This must be tracked
carefully in every formula that couples to the PPN metric.
The mapping \eqref{eq:sign-map} is the definitive translation.}


%%======================================================================
\section{L-1: Full 10-Parameter PPN Formalism}\label{sec:ppn}
%%======================================================================

\subsection{The Standard PPN Metric}

\fromlit{Will (2014), arXiv:1403.7377, \S3.2, Eq.~(36)}

The standard PPN metric in the standard PPN gauge, using
isotropic coordinates $(t, x^i)$ and the Newtonian potential
$U(\mathbf{x}) = \int \rho'(\mathbf{x}')\,|\mathbf{x} -
\mathbf{x}'|^{-1}\,\dd^3 x'$:

\begin{align}
  g_{00} &= -1 + 2U - 2\beta U^2
    - 2\xi\Phi_W
    + (2\gamma + 2 + \alpha_3 + \zeta_1 - 2\xi)\,\Phi_1
    \notag\\
    &\quad + 2(3\gamma - 2\beta + 1 + \zeta_2 + \xi)\,\Phi_2
    + 2(1 + \zeta_3)\,\Phi_3
    + 2(3\gamma + 3\zeta_4 - 2\xi)\,\Phi_4
    \notag\\
    &\quad - (\zeta_1 - 2\xi)\,\mathcal{A}
    - (\alpha_1 - \alpha_2 - \alpha_3)\,w^2 U
    - \alpha_2 w^i w^j U_{ij}
    \notag\\
    &\quad + (2\alpha_3 - \alpha_1)\,w^i V_i
    + \mathcal{O}(4),
    \label{eq:ppn-g00}\\[8pt]
  g_{0i} &= -\tfrac{1}{2}(4\gamma + 3 + \alpha_1
    - \alpha_2 + \zeta_1 - 2\xi)\,V_i
    - \tfrac{1}{2}(1 + \alpha_2 - \zeta_1 + 2\xi)\,W_i
    \notag\\
    &\quad - \tfrac{1}{2}(\alpha_1 - 2\alpha_2)\,w^i U
    - \alpha_2 w^j U_{ij}
    + \mathcal{O}(3.5),
    \label{eq:ppn-g0i}\\[8pt]
  g_{ij} &= (1 + 2\gamma U)\,\delta_{ij}
    + \mathcal{O}(2),
    \label{eq:ppn-gij}
\end{align}
where $\mathcal{O}(n)$ denotes terms of $n$-th post-Newtonian order.
The \emph{superpotentials} $\Phi_1, \ldots, \Phi_4, \Phi_W,
\mathcal{A}$ and the vector potentials $V_i, W_i$ are defined in
terms of matter integrals (Will~2014, \S3.2, Eqs.~(28)--(35)).
For the static, single-body problem that is relevant for PPN-1,
only the leading terms survive:
\begin{equation}\label{eq:ppn-static}
  g_{00} \approx -1 + 2U - 2\beta U^2,\qquad
  g_{ij} \approx (1 + 2\gamma U)\,\delta_{ij}.
\end{equation}

\subsection{The 10 PPN Parameters}

\fromlit{Will (2014), arXiv:1403.7377, Table~2}

\begin{definition}[PPN parameters and their physical meaning]
\label{def:ppn-params}
The 10 PPN parameters and their physical content:
\begin{center}
\renewcommand{\arraystretch}{1.25}
\begin{longtable}{clp{7.5cm}c}
\toprule
\# & Parameter & Physical meaning & GR value \\
\midrule
\endhead
1 & $\gamma$ & Spatial curvature produced per unit rest mass.
  Measures how much space curvature is produced by gravity. & 1 \\
2 & $\beta$ & Nonlinearity of gravitational superposition.
  Measures nonlinear contribution to $g_{00}$. & 1 \\
3 & $\xi$ & Preferred-location effects (Whitehead-type). & 0 \\
4 & $\alpha_1$ & Preferred-frame effect: PPN--velocity-dependent
  contribution to $g_{00}$ proportional to $w^2 U$. & 0 \\
5 & $\alpha_2$ & Preferred-frame effect: contribution to $g_{00}$
  proportional to $w^i w^j U_{ij}$. & 0 \\
6 & $\alpha_3$ & Preferred-frame effect combined with
  conservation-law violation.  & 0 \\
7 & $\zeta_1$ & Violation of conservation of total
  4-momentum (matter creation/destruction). & 0 \\
8 & $\zeta_2$ & Violation of conservation of total
  4-momentum. & 0 \\
9 & $\zeta_3$ & Violation of conservation of total
  4-momentum. & 0 \\
10 & $\zeta_4$ & Violation of conservation of total
  4-momentum. & 0 \\
\bottomrule
\end{longtable}
\end{center}
\end{definition}

\begin{remark}[Fully conservative theories]
\fromlit{Will (2014), \S3.3}
A metric theory of gravity derived from a generally covariant
Lagrangian automatically satisfies:
\begin{equation}\label{eq:conservative}
  \alpha_3 = \zeta_1 = \zeta_2 = \zeta_3 = \zeta_4 = 0.
\end{equation}
If additionally the theory has no preferred-frame structure
(no prior geometry, no fixed timelike vector), then:
\begin{equation}\label{eq:no-preferred}
  \xi = \alpha_1 = \alpha_2 = 0.
\end{equation}
A theory satisfying both conditions is called \emph{fully
conservative} and is described by only 2 independent PPN parameters:
$\gamma$ and $\beta$.
\end{remark}

\begin{finding}[SCT is fully conservative]\label{find:sct-conservative}
The SCT spectral action $\Tr(f(D^2/\Lambda^2))$ is derived from
a diffeomorphism-invariant functional with no prior geometry.
The off-shell gauge invariance proven in NT-4a (\S3) is the
linearized manifestation of this symmetry.  Therefore SCT is fully
conservative:
\begin{equation}\label{eq:sct-conservative}
  \xi = \alpha_1 = \alpha_2 = \alpha_3
  = \zeta_1 = \zeta_2 = \zeta_3 = \zeta_4 = 0.
\end{equation}
Only $\gamma$ and $\beta$ are nontrivial.  This is \textbf{derived},
not asserted.
\end{finding}

\subsection{The Nordtvedt Effect}

\fromlit{Will (2014), arXiv:1403.7377, \S3.7, Eq.~(56)}

The Nordtvedt effect quantifies the violation of the strong
equivalence principle (SEP) for gravitationally bound bodies.
The Nordtvedt parameter is:
\begin{equation}\label{eq:nordtvedt}
  \eta_N = 4\beta - \gamma - 3
    - \frac{10}{3}\xi - \alpha_1
    + \frac{2}{3}\alpha_2 - \frac{2}{3}\zeta_1
    - \frac{1}{3}\zeta_2.
\end{equation}
For a fully conservative theory (all $\xi, \alpha_i, \zeta_i = 0$):
\begin{equation}\label{eq:nordtvedt-simple}
  \eta_N = 4\beta - \gamma - 3.
\end{equation}
In GR ($\gamma = \beta = 1$): $\eta_N = 4 - 1 - 3 = 0$.

\signcorrection{The Nordtvedt parameter $\eta_N$ requires
knowledge of \emph{both} $\gamma$ and $\beta$.  The PPN-1
derivation computes $\gamma$ from the linearized propagator,
but $\beta$ requires the cubic vertex (nonlinear interaction).
The Nordtvedt effect bounds are therefore \emph{informational
only} for the linear PPN-1 analysis; they become constraining
once $\beta$ is derived from NT-4b.}


%%======================================================================
\section{L-2: Stelle (1977/78) Higher-Derivative Gravity}
\label{sec:stelle}
%%======================================================================

\subsection{The Quadratic Gravity Action}

\fromlit{Stelle, Phys.\ Rev.\ D \textbf{16} (1977) 953, Eq.~(1);
Gen.\ Rel.\ Grav.\ \textbf{9} (1978) 353}

The most general parity-even, quadratic-curvature extension of GR
in $d = 4$:
\begin{equation}\label{eq:stelle-action}
  S_{\mathrm{Stelle}} = \int \dd^4 x\,\sqrt{-g}\left[
    \frac{R}{16\pi G}
    + \alpha\, R^2
    + \beta_S\, C_{\mu\nu\rho\sigma}C^{\mu\nu\rho\sigma}
  \right],
\end{equation}
where $\alpha$ and $\beta_S$ are dimensionful coupling constants
with $[\alpha] = [\beta_S] = \mathrm{length}^2$.

\begin{convention}[Stelle coupling notation]\label{conv:stelle-couplings}
We use $\beta_S$ (with subscript) for the Stelle Weyl-squared
coupling to avoid confusion with the PPN parameter $\beta$.
Some references use the Ricci-squared basis:
$S = \int \dd^4 x\,\sqrt{-g}\,[R/(16\pi G)
+ c_1 R^2 + c_2 R_{\mu\nu}R^{\mu\nu}]$.
The conversion uses the 4D Gauss--Bonnet identity
$C_{\mu\nu\rho\sigma}C^{\mu\nu\rho\sigma}
= R_{\mu\nu\rho\sigma}R^{\mu\nu\rho\sigma}
- 2R_{\mu\nu}R^{\mu\nu} + \tfrac{1}{3}R^2$:
\begin{equation}\label{eq:stelle-basis}
  c_2 = 2\beta_S,\qquad
  c_1 = \alpha + \tfrac{1}{3}\beta_S.
\end{equation}
\end{convention}

\subsection{Linearized Propagator}

\fromlit{Stelle (1977), Eq.~(2.5); Rivers, Nuovo Cim.\ \textbf{34} (1964) 386}

The linearized graviton propagator around flat Minkowski space,
decomposed via Barnes--Rivers projectors:
\begin{equation}\label{eq:stelle-prop}
  G_{\mu\nu\rho\sigma}^{\mathrm{Stelle}}(k)
  = \frac{1}{k^2}\left[
    \frac{P^{(2)}_{\mu\nu\rho\sigma}}{1 + 2\beta_S\, k^2}
    + \frac{P^{(0\text{-}s)}_{\mu\nu\rho\sigma}}
           {1 + (6\alpha + 2\beta_S)\,k^2}
  \right]
  + \text{gauge terms},
\end{equation}
where the gauge-dependent terms drop out for on-shell matrix elements
with conserved sources.

\begin{definition}[Stelle effective masses]\label{def:stelle-masses}
\fromlit{Stelle (1977), below Eq.~(2.5)}
\begin{equation}\label{eq:stelle-masses}
  m_2^2 = \frac{1}{2\beta_S} \quad\text{(spin-2 massive mode)},\qquad
  m_0^2 = \frac{1}{6\alpha + 2\beta_S} \quad\text{(spin-0 massive mode)}.
\end{equation}
The spin-2 mass requires $\beta_S > 0$ for a real mass.
The spin-0 mass requires $3\alpha + \beta_S > 0$.
\end{definition}

\subsection{Partial-Fraction Decomposition}

The propagator denominators factor as:
\begin{align}
  \frac{1}{k^2(1 + k^2/m_2^2)}
  &= \frac{1}{k^2} - \frac{1}{k^2 + m_2^2}
  \quad\text{(spin-2)},
  \label{eq:pf-spin2}\\
  \frac{1}{k^2(1 + k^2/m_0^2)}
  &= \frac{1}{k^2} - \frac{1}{k^2 + m_0^2}
  \quad\text{(spin-0)}.
  \label{eq:pf-spin0}
\end{align}
The \emph{relative sign} between the two terms in \eqref{eq:pf-spin2}
is the signature of the ghost: the massive spin-2 propagator has
residue $-1$ (wrong sign), meaning it propagates a ghost degree of freedom.
The massive spin-0 has residue $+1$ (healthy).

\subsection{The Newtonian Potential}

\fromlit{Stelle (1978), Gen.\ Rel.\ Grav.\ \textbf{9}, 353;
Accioly et al.\ (2002), Phys.\ Rev.\ D \textbf{65}, 067501}

Performing the Fourier--Bessel integral for a static point mass $M$:

\begin{benchmark}[Stelle metric potentials]\label{bench:stelle-pots}
\begin{align}
  \Phi_{\mathrm{Stelle}}(r)
  &= -\frac{GM}{r}\left[
    1 - \frac{4}{3}\,e^{-m_2 r} + \frac{1}{3}\,e^{-m_0 r}
  \right],
  \label{eq:Phi-Stelle}\\
  \Psi_{\mathrm{Stelle}}(r)
  &= -\frac{GM}{r}\left[
    1 - \frac{2}{3}\,e^{-m_2 r} - \frac{1}{3}\,e^{-m_0 r}
  \right].
  \label{eq:Psi-Stelle}
\end{align}
\end{benchmark}

\noindent The coefficient structure is:
\begin{center}
\begin{tabular}{lccc}
\toprule
Mode & $\Phi$ coefficient & $\Psi$ coefficient & Physical character \\
\midrule
Massless (GR) & $+1$ & $+1$ & Healthy \\
Massive spin-2 ($m_2$) & $-4/3$ & $-2/3$ & Ghost (attractive Yukawa) \\
Massive spin-0 ($m_0$) & $+1/3$ & $-1/3$ & Healthy (repulsive Yukawa) \\
\midrule
Sum at $r = 0$ & $1 - 4/3 + 1/3 = 0$ & $1 - 2/3 - 1/3 = 0$ & \\
\bottomrule
\end{tabular}
\end{center}

\begin{finding}[Regularity at the origin]\label{find:V0}
$\Phi(0) = \Psi(0) = 0$: the potential is \emph{finite} (in fact,
vanishing) at $r = 0$ because the coefficients sum to zero:
$1 - 4/3 + 1/3 = 0$.  This is a direct consequence of the improved
UV behavior of the higher-derivative propagator.
\end{finding}

\begin{finding}[The ghost problem]\label{find:ghost}
The spin-2 massive mode has the \emph{wrong-sign} residue in the
propagator (negative residue in the spin-2 channel, Eq.~\eqref{eq:pf-spin2}).
This means:
\begin{itemize}[nosep]
  \item The massive spin-2 mode is a ghost (negative norm state).
  \item Its Yukawa contribution to $\Phi$ is \emph{attractive}
    (coefficient $-4/3 < 0$, enhancing the gravitational potential).
  \item The coefficient $-4/3$ vs $+1/3$ ensures $|$ghost$|
    > |$scalar$|$, so the total potential is enhanced at short
    distances before the $r \to 0$ cancellation.
\end{itemize}
The SCT treatment of this ghost pole is via the fakeon prescription
(Anselmi 2017, arXiv:1704.07728); see MR-2 for the full analysis.
\end{finding}

\subsection{Mapping to SCT Coefficients}

The SCT effective action in its local limit maps to Stelle's
framework via \fromlit{NT-1b Phase~3; NT-4a}:
\begin{equation}\label{eq:sct-to-stelle}
  \beta_S \leftrightarrow \frac{\alpha_C}{16\pi^2\,\Lambda^2}
  = \frac{13}{1920\pi^2\,\Lambda^2},
  \qquad
  \alpha \leftrightarrow \frac{\alpha_R(\xi) - \alpha_C/3}
  {16\pi^2\,\Lambda^2}.
\end{equation}
The full correspondence for the propagator denominators:
\begin{equation}\label{eq:Pi-stelle-map}
  \Pi_{\mathrm{TT}}^{\mathrm{Stelle}} = 1 + 2\beta_S\, k^2
  \longleftrightarrow 1 + \frac{13}{60}\,z
  \quad (z = k^2/\Lambda^2),
\end{equation}
\begin{equation}
  \Pi_{\mathrm{s}}^{\mathrm{Stelle}} = 1 + (6\alpha + 2\beta_S)\,k^2
  \longleftrightarrow 1 + 6\left(\xi - \tfrac{1}{6}\right)^2 z.
\end{equation}

The SCT effective masses in the local (Yukawa) approximation:
\begin{equation}\label{eq:sct-masses}
  m_2^2 = \frac{60}{13}\,\Lambda^2 \approx (2.148\,\Lambda)^2,
  \qquad
  m_0^2 = \frac{\Lambda^2}{6(\xi - 1/6)^2},
\end{equation}
with $m_0 \to \infty$ at conformal coupling $\xi = 1/6$
(scalar mode decouples).


%%======================================================================
\section{L-3: Nonlocal Gravity Potentials}\label{sec:nonlocal}
%%======================================================================

\subsection{The BMS Infinite-Derivative Action}

\fromlit{Biswas, Mazumdar, Siegel (2006), JCAP \textbf{0603}, 009;
hep-th/0508194, Eq.~(1)}

The most general quadratic-curvature action with nonlocal
(infinite-derivative) form factors:
\begin{equation}\label{eq:idg-action}
  S_{\mathrm{IDG}} = \frac{1}{16\pi G}\int \dd^4 x\,\sqrt{-g}\left[
    R + R\,\mathcal{F}_1(\Box/M^2)\,R
    + R_{\mu\nu}\,\mathcal{F}_2(\Box/M^2)\,R^{\mu\nu}
    + C_{\mu\nu\rho\sigma}\,\mathcal{F}_3(\Box/M^2)\,
    C^{\mu\nu\rho\sigma}
  \right],
\end{equation}
where $\mathcal{F}_i$ are analytic functions of the covariant
d'Alembertian $\Box = g^{\mu\nu}\nabla_\mu\nabla_\nu$, and $M$ is
the scale of non-locality.

\begin{remark}[Comparison with SCT action]
The SCT effective action has the same quadratic-curvature structure
but uses only $\Weylsq$ and $R^2$ (not $\Ricsq$ independently),
because the Gauss--Bonnet term is topological in 4D.  The SCT form
factors $F_1(z)$ and $F_2(z,\xi)$ are not free but are
\emph{computed} from the spectral action.
\end{remark}

\subsection{The Modified Propagator Structure}

\fromlit{Biswas, Gerwick, Koivisto, Mazumdar (2012), PRL \textbf{108},
031101; Edholm, Koshelev, Mazumdar (2016), arXiv:1604.01989, Eq.~(5)}

The linearized propagator around flat space decomposes as:
\begin{equation}\label{eq:idg-prop}
  G_{\mu\nu\rho\sigma}^{\mathrm{IDG}}(k)
  = \frac{P^{(2)}_{\mu\nu\rho\sigma}}{k^2\,a(-k^2/M^2)}
  + \frac{P^{(0\text{-}s)}_{\mu\nu\rho\sigma}}
         {k^2\,(a(-k^2/M^2) - 3c(-k^2/M^2))}
  + \text{gauge},
\end{equation}
where the two structure functions are
\fromlit{Biswas et al.\ (2012), Eqs.~(4)--(5)}:
\begin{align}
  a(-k^2/M^2) &= 1 - \frac{2}{M_{\mathrm{Pl}}^2}\,k^2\,
    \big(\mathcal{F}_2(-k^2/M^2) + 2\mathcal{F}_3(-k^2/M^2)\big),
    \label{eq:a-func}\\
  c(-k^2/M^2) &= 1 + \frac{2}{M_{\mathrm{Pl}}^2}\,k^2\,
    \big(6\mathcal{F}_1(-k^2/M^2)
    + \tfrac{1}{2}\mathcal{F}_2(-k^2/M^2)\big).
    \label{eq:c-func}
\end{align}

The mapping to the SCT notation:
\begin{equation}\label{eq:idg-to-sct}
  a(-k^2/M^2) \longleftrightarrow \Pi_{\mathrm{TT}}(z),\qquad
  a(-k^2/M^2) - 3c(-k^2/M^2) \longleftrightarrow \Pi_{\mathrm{s}}(z,\xi).
\end{equation}

\subsection{Ghost-Freedom Condition}

\fromlit{BMS (2006), hep-th/0508194, \S3; Biswas et al.\ (2012),
PRL \textbf{108}, 031101}

The propagator \eqref{eq:idg-prop} is ghost-free if and only if
\emph{neither} $a(-k^2/M^2)$ nor $(a - 3c)(-k^2/M^2)$ has zeros
on the real $k^2$ axis.  A sufficient condition is:
\begin{equation}\label{eq:ghost-free}
  a(\Box/M^2) = c(\Box/M^2) = e^{-\tau(\Box/M^2)},
\end{equation}
where $\tau$ is an \emph{entire function} with $\tau(0) = 0$
and $\operatorname{Re}\tau(-k^2/M^2) \geq 0$ for real $k^2 \geq 0$.

\begin{remark}[SCT does NOT satisfy the ghost-free condition]
\label{rem:sct-ghost}
In SCT, $\Pi_{\mathrm{TT}}(z)$ has a real positive zero at
$z_0 \approx 2.415$ (verified numerically in NT-4a).
Therefore the SCT spin-2 sector contains a massive ghost pole.
The SCT resolution is the fakeon prescription, not the
exponential suppression used in IDG.
\end{remark}

\subsection{The Edholm--Koshelev--Mazumdar Newtonian Potential}

\fromlit{Edholm, Koshelev, Mazumdar (2016), arXiv:1604.01989,
Eqs.~(8)--(10)}

For the simplest ghost-free choice $\tau = -k^2/M^2$
(Gaussian form factor), both $a = c = e^{k^2/M^2}$, and
the Newtonian potential becomes:
\begin{equation}\label{eq:Phi-idg}
  \Phi_{\mathrm{IDG}}(r) = -\frac{GM}{r}\,\erf\!\left(\frac{Mr}{2}\right),
\end{equation}
where $\erf(x) = \frac{2}{\sqrt{\pi}}\int_0^x e^{-t^2}\,\dd t$ is
the error function.

More generally, the exact potential is given by the Fourier--Bessel
integral \fromlit{EKM (2016), Eq.~(7)}:
\begin{equation}\label{eq:idg-FB}
  \Phi(r) = -\frac{GM}{r}\,\frac{2}{\pi}\int_0^\infty
  \frac{\sin(kr)}{k}\,K_\Phi(k^2/M^2)\,\dd k,
\end{equation}
where the Newton kernel $K_\Phi$ is constructed from the propagator
as described in \S\ref{sec:proj-pot}.

\begin{finding}[IDG forces $\gamma = 1$]\label{find:idg-gamma}
When $a = c$ (ghost-free condition), the spin-0 propagator denominator
is $a - 3c = -2a$, which is a \emph{fixed multiple} of the spin-2
denominator.  Consequently both metric potentials receive identical
modifications and $\gamma_{\mathrm{PPN}} = 1$ identically
\fromlit{EKM (2016), below Eq.~(10)}.
This is \emph{not} the case in SCT, where $\Pi_{\mathrm{TT}}
\neq \Pi_{\mathrm{s}}$ generically.
\end{finding}


%%======================================================================
\section{L-4: Brans--Dicke as Limiting Case}\label{sec:BD}
%%======================================================================

\subsection{The Brans--Dicke Action}

\fromlit{Brans, Dicke (1961), Phys.\ Rev.\ \textbf{124}, 925, Eq.~(5);
Will (2014), arXiv:1403.7377, \S5.1}

\begin{equation}\label{eq:BD-action}
  S_{\mathrm{BD}} = \frac{1}{16\pi}\int \dd^4 x\,\sqrt{-g}\left[
    \phi\, R - \frac{\omega}{\phi}\,
    g^{\mu\nu}\partial_\mu\phi\,\partial_\nu\phi
  \right] + S_{\mathrm{matter}},
\end{equation}
where $\phi$ is the Brans--Dicke scalar field and $\omega$ is the
dimensionless BD coupling constant.

\subsection{PPN Parameters of Brans--Dicke Theory}

\fromlit{Will (2014), arXiv:1403.7377, \S5.1, Eqs.~(80)--(82)}

\begin{benchmark}[Brans--Dicke PPN values]\label{bench:BD-ppn}
\begin{equation}\label{eq:gamma-BD}
  \gamma_{\mathrm{BD}} = \frac{1 + \omega}{2 + \omega},
  \qquad
  \beta_{\mathrm{BD}} = 1,
  \qquad
  \xi = \alpha_1 = \alpha_2 = \alpha_3 = \zeta_i = 0.
\end{equation}
\end{benchmark}

\noindent Key limits:
\begin{itemize}[nosep]
  \item GR limit: $\omega \to \infty$ gives $\gamma_{\mathrm{BD}} \to 1$.
  \item The deviation: $\gamma_{\mathrm{BD}} - 1
    = -1/(2 + \omega) \approx -1/\omega$ for large $\omega$.
  \item Cassini bound $|\gamma - 1| < 2.3\times 10^{-5}$ translates to
    $\omega > 40{,}000$.
\end{itemize}

\subsection{Structural Analogy with SCT}

\begin{finding}[Key structural difference from Brans--Dicke]
\label{find:bd-vs-sct}
In Brans--Dicke theory, $\gamma - 1$ is a \emph{constant}
(scale-independent) offset from GR.  This is fundamentally different
from SCT, where:
\begin{enumerate}[nosep]
  \item The scalar sector is controlled by $6(\xi - 1/6)^2$:
    at conformal coupling $\xi = 1/6$, the scalar mode decouples
    \emph{identically} and $\gamma = 1$ exactly, analogous to
    $\omega \to \infty$.
  \item For $\xi \neq 1/6$, the scalar mode is active but its
    contribution is \emph{Yukawa-suppressed}: $\Delta\gamma \sim
    e^{-m_0 r}$, not constant.
  \item At solar-system distances, $\Delta\gamma \sim e^{-10^{46}}
    \approx 0$ for $\Lambda \sim M_{\mathrm{Pl}}$.
\end{enumerate}
Therefore SCT evades the Cassini bound \emph{automatically} at any
$\xi$ (due to the Yukawa suppression), while Brans--Dicke requires
fine-tuning $\omega > 40{,}000$.
\end{finding}

\begin{finding}[BD as a consistency check]\label{find:bd-check}
The BD theory provides a structural consistency check for the SCT
$\gamma$ formula: in the limit where the SCT scalar mass $m_0 \to 0$
(a hypothetical limit, not realized in the actual theory), the
$\gamma$ formula should exhibit a BD-like constant offset from 1.
Specifically, the Yukawa approximation gives $\gamma(r) - 1 \sim
(2/3)\,e^{-m_0 r}$ for the scalar contribution.  As $m_0 r \to 0$:
$\gamma - 1 \to 2/3$, corresponding to an effective $\omega = -1/2$
(i.e., the massless limit is not BD but rather a conformal scalar
with $\gamma = 1/2$).  This matches the vDVZ value for a massless
scalar mediator, providing an internal consistency check.
\end{finding}


%%======================================================================
\section{L-5: Barnes--Rivers Projector Algebra}\label{sec:projectors}
%%======================================================================

\subsection{Projector Definitions}

\fromlit{van Nieuwenhuizen (1973), Nucl.\ Phys.\ B \textbf{60}, 478;
Rivers (1964), Nuovo Cim.\ \textbf{34}, 386;
NT-4a, Eqs.~(4.7)--(4.8)}

\begin{definition}[Transverse projector]\label{def:theta}
In flat-space momentum representation:
\begin{equation}\label{eq:theta}
  \theta_{\mu\nu}(k) = \delta_{\mu\nu} - \frac{k_\mu k_\nu}{k^2},
  \qquad
  \omega_{\mu\nu}(k) = \frac{k_\mu k_\nu}{k^2},
\end{equation}
satisfying $\theta_{\mu\nu} + \omega_{\mu\nu} = \delta_{\mu\nu}$,
$\theta_{\mu\alpha}\theta^{\alpha}{}_\nu = \theta_{\mu\nu}$,
$\theta^\mu{}_\mu = d - 1$ (in $d$ dimensions).
\end{definition}

\begin{definition}[Spin-2 and spin-0 projectors]\label{def:BR}
\begin{align}
  P^{(2)}_{\mu\nu\rho\sigma}
  &= \frac{1}{2}(\theta_{\mu\rho}\theta_{\nu\sigma}
  + \theta_{\mu\sigma}\theta_{\nu\rho})
  - \frac{1}{d-1}\,\theta_{\mu\nu}\theta_{\rho\sigma},
  \label{eq:P2}\\
  P^{(0\text{-}s)}_{\mu\nu\rho\sigma}
  &= \frac{1}{d-1}\,\theta_{\mu\nu}\theta_{\rho\sigma}.
  \label{eq:P0s}
\end{align}
In $d = 4$: the prefactor $1/(d-1) = 1/3$.
\end{definition}

\subsection{Key Algebraic Properties}

\fromlit{van Nieuwenhuizen (1973), \S2}

\begin{enumerate}[nosep,label=(\roman*)]
  \item \textbf{Idempotency:}
    $P^{(J)}\cdot P^{(J)} = P^{(J)}$ for $J = 2, 0$.
  \item \textbf{Orthogonality:}
    $P^{(2)}\cdot P^{(0\text{-}s)} = 0$.
  \item \textbf{Traces:}
    \begin{equation}\label{eq:traces}
      \tr\!\left(P^{(2)}\right)
      \equiv P^{(2)\,\mu\nu}{}_{\mu\nu}
      = \frac{d(d-1)}{2} - 1
      = 5 \quad (d=4),
    \end{equation}
    \begin{equation}
      \tr\!\left(P^{(0\text{-}s)}\right)
      = 1.
    \end{equation}
  \item \textbf{Partial completeness:}
    $P^{(2)} + P^{(0\text{-}s)} = P^T$, where $P^T$ is the
    transverse projector:
    \begin{equation}
      P^{T}_{\mu\nu\rho\sigma}
      = \frac{1}{2}(\theta_{\mu\rho}\theta_{\nu\sigma}
      + \theta_{\mu\sigma}\theta_{\nu\rho}).
    \end{equation}
  \item \textbf{Full decomposition} (van~Nieuwenhuizen):
    \begin{equation}\label{eq:full-decomp}
      \delta_{(\mu}{}^{(\rho}\,\delta_{\nu)}{}^{\sigma)}
      = P^{(2)} + P^{(0\text{-}s)} + P^{(0\text{-}w)}
      + P^{(1)} + P^{(0\text{-}sw)} + P^{(0\text{-}ws)},
    \end{equation}
    where $P^{(1)}$, $P^{(0\text{-}w)}$, and the cross terms
    involve the longitudinal projector $\omega_{\mu\nu}$ and are
    pure gauge.
\end{enumerate}

\subsection{Projector Values for Static Point Source}

\begin{finding}[Critical for PPN-1 derivation]\label{find:proj-static}
For a static point mass with $k_0 = 0$ (purely spatial momentum
$\vec{k}$), the transverse projector evaluates as:
\begin{equation}\label{eq:theta-static}
  \theta_{\mu\nu}\big|_{k_0=0}
  = \begin{cases}
    1 & \mu = \nu = 0,\\
    \delta_{ij} - \hat{k}_i\hat{k}_j & \mu = i,\, \nu = j,\\
    0 & \text{mixed } (\mu = 0, \nu = i \text{ or vice versa}).
  \end{cases}
\end{equation}
The spatial trace is
$\theta_{ii}\big|_{k_0=0} = 3 - 1 = 2$ (in $d = 3$ spatial
dimensions), and the full trace is
$\theta_{\mu\mu}\big|_{k_0=0} = 1 + 2 = 3$.
\end{finding}

The projector components relevant for the $h_{00}$ and $h_{ij}$
extraction:
\begin{align}
  P^{(2)}_{0000} &= \theta_{00}\theta_{00}
    - \tfrac{1}{3}\theta_{00}\theta_{00}
    = 1 - \tfrac{1}{3} = \tfrac{2}{3},
    \label{eq:P2-0000}\\
  P^{(0\text{-}s)}_{0000} &= \tfrac{1}{3}\theta_{00}\theta_{00}
    = \tfrac{1}{3},
    \label{eq:P0s-0000}\\
  P^{(2)}_{00ij} &= -\tfrac{1}{3}\theta_{ij},
    \label{eq:P2-00ij}\\
  P^{(0\text{-}s)}_{00ij} &= \tfrac{1}{3}\theta_{ij}.
    \label{eq:P0s-00ij}
\end{align}
These are the values that enter the projector-to-potential map
(see \S\ref{sec:proj-pot}).

\signcorrection{In the previous PPN-1 attempt, the values
$P^{(2)}_{0000} = 2/3$ and $P^{(0\text{-}s)}_{0000} = 1/3$ were
extracted incorrectly, leading to a vanishing $h_{00}$ in the GR
limit (Finding~B).  The values above are verified against three
independent references: NT-4a, van~Nieuwenhuizen~(1973), and
direct computation with $\theta_{00} = 1$.}


%%======================================================================
\section{L-5 (continued): The Projector-to-Potential Map}
\label{sec:proj-pot}
%%======================================================================

This section contains the \textbf{central technical result} of the
literature extraction: the explicit formulas relating the propagator
components $\Pi_{\mathrm{TT}}$ and $\Pi_{\mathrm{s}}$ to the two
metric potentials $\Phi(r)$ and $\Psi(r)$.

\subsection{Statement of the Map}

\begin{theorem}[Projector-to-potential map]\label{thm:proj-pot}
For a modified graviton propagator of the form
\begin{equation}\label{eq:G-full}
  G_{\mu\nu\rho\sigma}(k)
  = \frac{1}{k^2}\left[
    \frac{P^{(2)}_{\mu\nu\rho\sigma}}{\Pi_{\mathrm{TT}}(z)}
    + \frac{P^{(0\text{-}s)}_{\mu\nu\rho\sigma}}{\Pi_{\mathrm{s}}(z,\xi)}
  \right],
\end{equation}
the static metric potentials for a point mass $M$ are:
\begin{align}
  \Phi(r) &= -\frac{GM}{r}\,\frac{2}{\pi}
  \int_0^\infty \frac{\sin(kr)}{k}\, K_\Phi(k^2/\Lambda^2)\,\dd k,
  \label{eq:Phi-exact}\\
  \Psi(r) &= -\frac{GM}{r}\,\frac{2}{\pi}
  \int_0^\infty \frac{\sin(kr)}{k}\, K_\Psi(k^2/\Lambda^2)\,\dd k,
  \label{eq:Psi-exact}
\end{align}
where the \emph{Newton kernels} are:
\begin{equation}\label{eq:K-Phi}
  \boxed{K_\Phi(z)
  = \frac{4}{3\,\Pi_{\mathrm{TT}}(z)}
  - \frac{1}{3\,\Pi_{\mathrm{s}}(z,\xi)},}
\end{equation}
\begin{equation}\label{eq:K-Psi}
  \boxed{K_\Psi(z)
  = \frac{2}{3\,\Pi_{\mathrm{TT}}(z)}
  + \frac{1}{3\,\Pi_{\mathrm{s}}(z,\xi)}.}
\end{equation}
\end{theorem}

\subsection{Derivation Sketch}

The derivation proceeds by:
\begin{enumerate}[nosep]
  \item Coupling the propagator to a static point-mass source
    $T_{\mu\nu}(k) = M\,\delta_{\mu 0}\delta_{\nu 0}$.
  \item Computing $\bar{h}_{\mu\nu}(k) = -(16\pi G/k^2)\,
    [P^{(2)}/\Pi_{\mathrm{TT}} + P^{(0\text{-}s)}/\Pi_{\mathrm{s}}]
    \cdot \bar{T}_{\rho\sigma}$ with trace-reversed source
    $\bar{T}_{\mu\nu} = T_{\mu\nu} - \frac{1}{2}\eta_{\mu\nu}T$.
  \item Evaluating the projector action using
    Eqs.~\eqref{eq:P2-0000}--\eqref{eq:P0s-00ij}.
  \item Trace-reversing $h_{\mu\nu} = \bar{h}_{\mu\nu}
    - \frac{1}{2}\delta_{\mu\nu}\bar{h}$.
  \item Performing the angular average over $\hat{k}$ to extract
    the isotropic (monopole) components.
  \item Identifying $h_{00}(r) = -2\Phi(r)$ and
    $h_{ij}(r) = -2\Psi(r)\,\delta_{ij}$.
\end{enumerate}
The full derivation is in \texttt{PPN1\_derivation.tex}.

\subsection{GR Limit Check}

At $z = 0$: $\Pi_{\mathrm{TT}}(0) = \Pi_{\mathrm{s}}(0) = 1$, so:
\begin{equation}
  K_\Phi(0) = \frac{4}{3} - \frac{1}{3} = 1,\qquad
  K_\Psi(0) = \frac{2}{3} + \frac{1}{3} = 1.
\end{equation}
Both potentials equal $-GM/r$: $\gamma_{\mathrm{PPN}} = 1$.

\verified{GR limit $K_\Phi(0) = K_\Psi(0) = 1$ confirmed analytically
and numerically (mpmath 100-digit precision).}

\subsection{Stelle Benchmark Check}

In the Stelle limit ($\Pi_{\mathrm{TT}} = 1 + k^2/m_2^2$,
$\Pi_{\mathrm{s}} = 1 + k^2/m_0^2$), using the standard Fourier--Bessel
identity
$\frac{2}{\pi}\int_0^\infty \frac{\sin(kr)}{k}\,
\frac{1}{1+k^2/m^2}\,\dd k = e^{-mr}$:
\begin{align}
  \Phi(r) &= -\frac{GM}{r}\left[
    \frac{4}{3} - \frac{1}{3}
    - \frac{4}{3}\,e^{-m_2 r} + \frac{1}{3}\,e^{-m_0 r}
  \right] \notag\\
  &= -\frac{GM}{r}\left[
    1 - \frac{4}{3}\,e^{-m_2 r} + \frac{1}{3}\,e^{-m_0 r}
  \right],
\end{align}
which matches \eqref{eq:Phi-Stelle} exactly.  Similarly for $\Psi(r)$.

\crosscheck{Both $K_\Phi$ and $K_\Psi$ reproduce the known Stelle
potentials.  This is the primary verification of the
projector-to-potential map.}

\subsection{Universal Coupling Ratios}

The coefficient pattern in the Newton kernels follows from the
projector algebra alone, not from the specific form of
$\Pi_{\mathrm{TT}}$ or $\Pi_{\mathrm{s}}$:

\begin{center}
\begin{tabular}{lccc}
\toprule
Mode & $K_\Phi$ weight & $K_\Psi$ weight & $K_\Psi/K_\Phi$ ratio \\
\midrule
Spin-2 ($\Pi_{\mathrm{TT}}$) & $+4/3$ & $+2/3$ & $+1/2$ \\
Spin-0 ($\Pi_{\mathrm{s}}$) & $-1/3$ & $+1/3$ & $-1$ \\
\bottomrule
\end{tabular}
\end{center}

These ratios are \emph{model-independent}: any theory with the
propagator structure \eqref{eq:G-full} will have the same coefficients.

\subsection{The Effective PPN Gamma}

\begin{definition}[Scale-dependent PPN gamma]\label{def:gamma-eff}
\begin{equation}\label{eq:gamma-eff}
  \boxed{\gamma_{\mathrm{eff}}(z)
  = \frac{K_\Psi(z)}{K_\Phi(z)}
  = \frac{2/(3\Pi_{\mathrm{TT}}) + 1/(3\Pi_{\mathrm{s}})}
         {4/(3\Pi_{\mathrm{TT}}) - 1/(3\Pi_{\mathrm{s}})}.}
\end{equation}
\end{definition}

\noindent Limiting cases:
\begin{enumerate}[nosep]
  \item \textbf{GR} ($z \to 0$):
    $\gamma_{\mathrm{eff}} = (2/3 + 1/3)/(4/3 - 1/3) = 1$.
  \item \textbf{Conformal coupling} ($\xi = 1/6$, $\Pi_{\mathrm{s}} = 1$):
    $\gamma_{\mathrm{eff}} = (2/(3\Pi_{\mathrm{TT}}) + 1/3)
    / (4/(3\Pi_{\mathrm{TT}}) - 1/3)$.
  \item \textbf{Large $r$} (Yukawa regime):
    all exponential corrections vanish, $\gamma \to 1$.
  \item \textbf{Yukawa approximation}:
    \begin{equation}\label{eq:gamma-yukawa}
      \gamma_{\mathrm{eff}}(r)
      = \frac{1 - \frac{2}{3}\,e^{-m_2 r}
        - \frac{1}{3}\,e^{-m_0 r}}
             {1 - \frac{4}{3}\,e^{-m_2 r}
        + \frac{1}{3}\,e^{-m_0 r}}.
    \end{equation}
\end{enumerate}


%%======================================================================
\section{L-6: Experimental Data Compilation}\label{sec:experiments}
%%======================================================================

\subsection{Solar System: Light Deflection and Time Delay}

\subsubsection{Cassini Spacecraft (2003)}

\fromlit{Bertotti, Iess, Tortora, Nature \textbf{425} (2003) 374}

The Cassini spacecraft measured the Shapiro time delay of radio
signals passing near the Sun during solar conjunction in June 2002.
The frequency shift of the radio link is proportional to $\gamma + 1$.

\begin{benchmark}[Cassini bound on $\gamma$]\label{bench:cassini}
\begin{equation}\label{eq:cassini}
  \gamma - 1 = (2.1 \pm 2.3) \times 10^{-5}
  \qquad\text{(68\% CL)}.
\end{equation}
This gives the $2\sigma$ bound:
\begin{equation}\label{eq:cassini-2sigma}
  |\gamma - 1| < 4.4 \times 10^{-5}
  \qquad\text{(95\% CL, Cassini)}.
\end{equation}
\end{benchmark}

\noindent The measurement probes the gravitational field at
$r \sim 1$--$10\,R_\odot \approx (7$--$70) \times 10^8$~m, where
$R_\odot = 6.957 \times 10^8$~m.

\subsubsection{MESSENGER (2014)}

\fromlit{Verma, Fienga, Laskar et al., A\&A \textbf{561} (2014) A115}

The MESSENGER spacecraft orbiting Mercury provided an independent
measurement of $\gamma$ via Mercury ranging data:

\begin{benchmark}[MESSENGER bound on $\gamma$]\label{bench:messenger}
\begin{equation}\label{eq:messenger}
  |\gamma - 1| < 2.5 \times 10^{-5}
  \qquad\text{(MESSENGER + planetary ephemeris, 1$\sigma$)}.
\end{equation}
\end{benchmark}

\noindent This is comparable in precision to Cassini and provides
an independent confirmation at a different orbital geometry.

\subsubsection{Very Long Baseline Interferometry (VLBI)}

\fromlit{Lambert, Le~Poncin-Lafitte, A\&A \textbf{499} (2009) 331}

VLBI measurements of radio quasar positions during solar conjunction:
\begin{equation}\label{eq:vlbi}
  \gamma - 1 = (-1.6 \pm 3.3) \times 10^{-4}
  \qquad\text{(VLBI, 1$\sigma$)}.
\end{equation}
Less constraining than Cassini but provides independent
geometric information.

\subsection{Nordtvedt Effect and Equivalence Principle}

\subsubsection{Lunar Laser Ranging (LLR)}

\fromlit{Williams, Turyshev, Boggs, PRL \textbf{93} (2004) 261101;
Williams, Turyshev, Boggs, CQG \textbf{29} (2012) 184004}

LLR measures the Earth--Moon distance with $\sim 1$~cm precision,
sensitive to the Nordtvedt effect (gravitational self-energy
contribution to free fall):

\begin{benchmark}[LLR bound on the Nordtvedt parameter]
\label{bench:llr}
\begin{equation}\label{eq:llr-eta}
  |\eta_N| < 4.4 \times 10^{-4}
  \qquad\text{(95\% CL, LLR 2004)}.
\end{equation}
Using $\eta_N = 4\beta - \gamma - 3$ (for conservative theories), and
assuming $\gamma = 1$ (Cassini):
\begin{equation}\label{eq:llr-beta}
  |\beta - 1| < 1.1 \times 10^{-4}
  \qquad\text{(derived from LLR + Cassini)}.
\end{equation}
\end{benchmark}

\signcorrection{The LLR bound on $|\beta - 1|$ requires the
\emph{combination} of LLR and Cassini data, not LLR alone.
Quoting $|\beta - 1| < 8 \times 10^{-5}$ (as in Will~2014, Table~4)
uses additional constraints from perihelion precession.}

\subsubsection{MICROSCOPE (2022)}

\fromlit{Touboul et al., PRL \textbf{129} (2022) 121102}

The MICROSCOPE satellite tested the weak equivalence principle (WEP)
in Earth orbit by comparing the free-fall acceleration of
titanium and platinum test masses:

\begin{benchmark}[MICROSCOPE bound on WEP violation]
\label{bench:microscope}
\begin{equation}\label{eq:microscope}
  \eta_{\mathrm{WEP}}(\mathrm{Ti,Pt})
  = [-1.5 \pm 2.3\,(\mathrm{stat})
    \pm 1.5\,(\mathrm{syst})] \times 10^{-15}.
\end{equation}
The combined $2\sigma$ bound:
\begin{equation}\label{eq:microscope-bound}
  |\eta_{\mathrm{WEP}}| < 1.5 \times 10^{-15}
  \qquad\text{(MICROSCOPE, $2\sigma$ upper limit)}.
\end{equation}
\end{benchmark}

\begin{remark}[WEP vs Nordtvedt effect]
MICROSCOPE tests the \emph{weak} equivalence principle (composition
dependence of free fall for laboratory-scale bodies), not the
\emph{strong} equivalence principle (gravitational self-energy
contribution).  The Nordtvedt effect is a SEP violation, tested
by LLR.  MICROSCOPE's result $|\eta| < 10^{-15}$ constrains
composition-dependent fifth forces, not the PPN parameter $\eta_N$
directly.  However, any Yukawa-type force from SCT is
\emph{composition-independent} (universal coupling to mass),
so MICROSCOPE provides no additional constraint beyond the ISL
tests at the relevant range.
\end{remark}

\subsection{Short-Range Gravity: Torsion Balance}

\subsubsection{E\"ot-Wash Group}

\fromlit{Lee, Adelberger, Cook et al., PRL \textbf{124} (2020) 101101;
Kapner et al., PRL \textbf{98} (2007) 021101}

The E\"ot-Wash torsion balance experiments test the gravitational
inverse-square law (ISL) at sub-millimeter distances by searching
for Yukawa deviations of the form:
\begin{equation}\label{eq:yukawa-param}
  V(r) = -\frac{Gm_1 m_2}{r}\left[
    1 + \alpha\,e^{-r/\lambda}
  \right],
\end{equation}
where $\alpha$ is the coupling strength relative to gravity and
$\lambda$ is the Yukawa range.

\begin{benchmark}[E\"ot-Wash ISL exclusion]\label{bench:eotwash}
The current state-of-the-art exclusion
\fromlit{Lee et al.\ (2020)}:
\begin{equation}\label{eq:eotwash}
  |\alpha| \leq 1 \quad\text{at}\quad
  \lambda \geq 38.6\;\mu\mathrm{m}
  \qquad\text{(95\% CL)}.
\end{equation}
Previous result \fromlit{Kapner et al.\ (2007)}:
$|\alpha| \leq 1$ at $\lambda \geq 56\;\mu$m.
\end{benchmark}

\noindent The exclusion curve is approximately:
\begin{equation}\label{eq:eotwash-curve}
  |\alpha(\lambda)|_{\text{95\%}} \approx
  \begin{cases}
    1 & \lambda \gtrsim 38.6\;\mu\text{m},\\
    (\lambda / 38.6\;\mu\text{m})^{-4} & \lambda \lesssim 38.6\;\mu\text{m}.
  \end{cases}
\end{equation}

The HUST group (Huazhong University) provides independent
confirmation \fromlit{Tan et al., PRL \textbf{124} (2020) 051301}:
$|\alpha| < 1$ at $\lambda \geq 40\;\mu$m.

\begin{finding}[Relevance for SCT]\label{find:eotwash-sct}
In SCT, the Yukawa deviations have:
\begin{itemize}[nosep]
  \item Spin-2: $\alpha_2 = -4/3$, $\lambda_2 = 1/m_2$.
  \item Spin-0 ($\xi \neq 1/6$): $\alpha_0 = +1/3$, $\lambda_0 = 1/m_0$.
\end{itemize}
The dominant mode for ISL tests is the spin-2 Yukawa with
$|\alpha_2| = 4/3 > 1$.  The E\"ot-Wash bound $|\alpha| \leq 1$
at $\lambda = 38.6\;\mu$m then requires:
\begin{equation}\label{eq:lambda-bound}
  \lambda_2 = \frac{1}{m_2} < 38.6\;\mu\mathrm{m}
  \quad\Longrightarrow\quad
  m_2 > \frac{1}{38.6\;\mu\mathrm{m}}
  \approx 5.1\;\mathrm{meV}.
\end{equation}
Since $m_2 = \sqrt{60/13}\,\Lambda \approx 2.15\,\Lambda$, this
gives:
\begin{equation}\label{eq:Lambda-bound}
  \Lambda > \frac{5.1\;\mathrm{meV}}{2.15}
  \approx 2.4\;\mathrm{meV}
  \qquad\text{(from E\"ot-Wash)}.
\end{equation}
A more precise analysis (accounting for the exact $\alpha$ value
$4/3 > 1$) is performed in LT-3d, yielding
$\Lambda > 2.565$~meV as the parameter-free lower bound.
\end{finding}

\subsection{Summary of Experimental Bounds}

\begin{center}
\renewcommand{\arraystretch}{1.2}
\begin{longtable}{p{3cm}p{4cm}p{3.5cm}p{1.2cm}}
\toprule
\textbf{Experiment} & \textbf{Reference} & \textbf{Bound} &
\textbf{Year} \\
\midrule
\endhead
Cassini &
  Bertotti, Iess, Tortora,
  Nature \textbf{425}, 374 &
  $|\gamma-1| = (2.1\pm 2.3)\times 10^{-5}$ &
  2003 \\
MESSENGER &
  Verma et al.,
  A\&A \textbf{561}, A115 &
  $|\gamma-1| < 2.5\times 10^{-5}$ &
  2014 \\
VLBI &
  Lambert, Le~Poncin-Lafitte,
  A\&A \textbf{499}, 331 &
  $|\gamma-1| < 3.4\times 10^{-4}$ &
  2009 \\
LLR &
  Williams et al.,
  PRL \textbf{93}, 261101 &
  $|\eta_N| < 4.4\times 10^{-4}$ &
  2004 \\
LLR + Cassini &
  Will (2014) &
  $|\beta-1| < 8 \times 10^{-5}$ &
  2014 \\
MICROSCOPE &
  Touboul et al.,
  PRL \textbf{129}, 121102 &
  $|\eta_{\mathrm{WEP}}| < 1.5\times 10^{-15}$ &
  2022 \\
E\"ot-Wash &
  Lee et al.,
  PRL \textbf{124}, 101101 &
  $|\alpha|<1$ at $\lambda = 38.6\;\mu$m &
  2020 \\
E\"ot-Wash &
  Kapner et al.,
  PRL \textbf{98}, 021101 &
  $|\alpha|<1$ at $\lambda = 56\;\mu$m &
  2007 \\
HUST &
  Tan et al.,
  PRL \textbf{124}, 051301 &
  $|\alpha|<1$ at $\lambda = 40\;\mu$m &
  2020 \\
\bottomrule
\end{longtable}
\end{center}


%%======================================================================
\section{L-7: The van~Dam--Veltman--Zakharov Discontinuity}
\label{sec:vdvz}
%%======================================================================

\subsection{Statement of the Problem}

\fromlit{van~Dam, Veltman, Nucl.\ Phys.\ B \textbf{22} (1970) 397;
Zakharov, JETP Lett.\ \textbf{12} (1970) 312}

\begin{theorem}[vDVZ discontinuity]\label{thm:vdvz}
A Fierz--Pauli massive spin-2 field $h_{\mu\nu}$ with mass $m$,
coupled to a conserved source via $h_{\mu\nu}T^{\mu\nu}$,
predicts a light-deflection parameter:
\begin{equation}\label{eq:vdvz-gamma}
  \gamma_{\mathrm{FP}} = \frac{1}{2}
  \qquad\text{for \emph{any} } m > 0.
\end{equation}
This value does \emph{not} approach the GR value $\gamma = 1$
as $m \to 0$.  The discontinuity arises because the Fierz--Pauli
mass term introduces 5 propagating degrees of freedom (the 5
polarizations of a massive spin-2) for any $m > 0$, while a
massless spin-2 has only 2.  The extra scalar polarization
contributes a finite amount to the static potential even as
$m \to 0$.
\end{theorem}

\begin{remark}[Physical consequence]
$\gamma = 1/2$ is excluded by the Cassini bound at
$> 10^4\sigma$.  Any theory with a Fierz--Pauli mass term
must invoke the Vainshtein mechanism (nonlinear screening)
or abandon the Fierz--Pauli structure.
\end{remark}

\subsection{Propagator Origin of the vDVZ Discontinuity}

\fromlit{van~Dam, Veltman (1970), \S3}

The massive Fierz--Pauli propagator in $d = 4$:
\begin{equation}\label{eq:FP-prop}
  G_{\mu\nu\rho\sigma}^{\mathrm{FP}}(k)
  = \frac{\bar{P}^{(2)}_{\mu\nu\rho\sigma}(k,m)}{k^2 + m^2}
  - \frac{\frac{1}{3}\bar{\theta}_{\mu\nu}\bar{\theta}_{\rho\sigma}}
         {k^2 + m^2},
\end{equation}
where $\bar{\theta}_{\mu\nu} = \eta_{\mu\nu} - k_\mu k_\nu/m^2$
and $\bar{P}^{(2)}$ is the massive spin-2 projector.  The
\emph{relative sign} between the spin-2 and spin-0 (trace) parts
is $-1/3$ (versus $+1/3$ in the massless case).  This sign
difference produces $\gamma = 1/2$ instead of $\gamma = 1$.

In the massless limit $m \to 0$, the massive projectors
$\bar{P}^{(J)}$ do not reduce to the massless projectors $P^{(J)}$:
the longitudinal polarizations survive and contribute at order $m^0$.

\subsection{Why SCT Avoids vDVZ}

\begin{finding}[SCT mechanism for avoiding vDVZ]\label{find:sct-vdvz}
The SCT spectral action avoids the vDVZ discontinuity for three
distinct reasons:
\begin{enumerate}[nosep]
  \item \textbf{No Fierz--Pauli mass term.}  The spin-2 mass in SCT
    ($m_2 \sim 2.15\,\Lambda$) arises from the momentum-dependent
    propagator modification $\Pi_{\mathrm{TT}}(z) = 1 + (13/60)\,z\,
    \hat{F}_1(z)$, not from a classical mass term $m^2 h_{\mu\nu}
    h^{\mu\nu}$.  The propagator structure is:
    \begin{equation}
      G_{\mathrm{TT}}(k) = \frac{1}{k^2\,\Pi_{\mathrm{TT}}(k^2/\Lambda^2)},
    \end{equation}
    which smoothly reduces to $1/k^2$ as $k \to 0$.  There is
    no discontinuity in the number of propagating degrees of
    freedom.

  \item \textbf{Gauge invariance is preserved.}  The SCT effective
    action is gauge-invariant (NT-4a, \S3), so the propagator is
    defined \emph{within the gauge-invariant sector} (transverse
    projectors $P^{(2)}$, $P^{(0\text{-}s)}$).  The
    Fierz--Pauli mass term breaks gauge invariance, introducing
    the Boulware--Deser ghost at nonlinear level.  SCT avoids
    this entirely.

  \item \textbf{The dressed propagator is smooth at $k = 0$.}
    $\Pi_{\mathrm{TT}}(0) = 1$ and $\Pi_{\mathrm{s}}(0) = 1$
    by construction (NT-4a, Theorem~5.1).  Therefore the
    propagator $1/(k^2\,\Pi)$ has a simple $1/k^2$ pole at
    $k = 0$, identical to GR.  The massive pole at
    $k^2 = -m_2^2 \Lambda^2$ (where $\Pi_{\mathrm{TT}}$
    vanishes) is a separate feature that does not affect the
    $k \to 0$ behavior.
\end{enumerate}
\end{finding}

\begin{remark}[Comparison with IDG]
Ghost-free infinite-derivative gravity (BMS, Modesto) also avoids
vDVZ by using entire-function form factors that modify the
propagator without introducing mass terms.  The IDG approach
\emph{additionally} avoids ghost poles, while SCT has ghost poles
that are treated via the fakeon prescription.  Both approaches
preserve gauge invariance and the smooth $k \to 0$ limit.
\end{remark}


%%======================================================================
\section{Convention Summary}\label{sec:conventions-summary}
%%======================================================================

\begin{center}
\renewcommand{\arraystretch}{1.25}
\begin{longtable}{p{3.5cm}p{7cm}p{2.5cm}}
\toprule
\textbf{Item} & \textbf{Convention} & \textbf{Source} \\
\midrule
\endhead
Metric signature &
  $(-,+,+,+)$ Lorentzian; Euclidean for derivation,
  continued at the end &
  NT-1 \\
Static metric &
  $\dd s^2 = -(1+2\Phi)\dd t^2
  + (1-2\Psi)\delta_{ij}\dd x^i\dd x^j$;
  $\Phi,\Psi < 0$ for attraction &
  NT-4a \\
Newtonian potential &
  $U = GM/r > 0$ (Will); $\Phi = -U < 0$ (ours) &
  Will (2014) / NT-4a \\
PPN $\gamma$ &
  $g_{ij} = (1+2\gamma U)\delta_{ij}$; equivalently
  $\gamma = \Psi/\Phi$ at Newtonian order &
  Will (2014) \\
PPN $\beta$ &
  $g_{00} = -1+2U-2\beta U^2 + \ldots$ &
  Will (2014) \\
Nordtvedt &
  $\eta_N = 4\beta-\gamma-3$ (conservative theories) &
  Will (2014) \\
Transverse projector &
  $\theta_{\mu\nu} = \delta_{\mu\nu} - k_\mu k_\nu/k^2$ &
  NT-4a \\
Barnes--Rivers &
  $P^{(2)}: 1/(d-1) = 1/3$ in $d=4$;
  $P^{(0\text{-}s)} = (1/3)\theta\theta$ &
  NT-4a \\
Projector traces &
  $\tr(P^{(2)}) = 5$, $\tr(P^{(0\text{-}s)}) = 1$ &
  van Nieuwenhuizen \\
Propagator &
  $G = k^{-2}[P^{(2)}/\Pi_{\mathrm{TT}} +
  P^{(0\text{-}s)}/\Pi_{\mathrm{s}}]$ &
  NT-4a \\
$\Pi_{\mathrm{TT}}$ &
  $1 + (13/60)\,z\,\hat{F}_1(z)$; $z = k^2/\Lambda^2$ &
  NT-4a Eq.~(5.3) \\
$\Pi_{\mathrm{s}}$ &
  $1 + 6(\xi-1/6)^2\,z\,\hat{F}_2(z,\xi)$ &
  NT-4a Eq.~(5.5) \\
$K_\Phi$ &
  $4/(3\Pi_{\mathrm{TT}}) - 1/(3\Pi_{\mathrm{s}})$ &
  This work \\
$K_\Psi$ &
  $2/(3\Pi_{\mathrm{TT}}) + 1/(3\Pi_{\mathrm{s}})$ &
  This work \\
Effective masses &
  $m_2^2 = 60\Lambda^2/13$;
  $m_0^2 = \Lambda^2/(6(\xi-1/6)^2)$ &
  NT-4a Eq.~(7.3) \\
SCT coefficients &
  $\alpha_C = 13/120$; $\alpha_R = 2(\xi-1/6)^2$;
  $c_2 = 13/60$ &
  NT-1b Phase~3 \\
Stelle couplings &
  $\beta_S = c_2/2 = 13/120$ (in units $1/(16\pi^2\Lambda^2)$);
  $c_1 = (\alpha_R - \alpha_C/3)/(16\pi^2)$ &
  Stelle (1977) \\
Preferred-frame &
  All zero: $\xi = \alpha_i = \zeta_i = 0$ (from gauge invariance) &
  Finding~\ref{find:sct-conservative} \\
Stelle potential &
  $V/V_N = 1 - (4/3)e^{-m_2 r} + (1/3)e^{-m_0 r}$ &
  Stelle (1978) \\
Yukawa parametrization &
  $V(r) = -(Gm_1 m_2/r)[1 + \alpha\,e^{-r/\lambda}]$ &
  Adelberger et al. \\
E\"ot-Wash ISL &
  $|\alpha| \leq 1$ at $\lambda = 38.6\;\mu$m &
  Lee et al.\ (2020) \\
\bottomrule
\end{longtable}
\end{center}


%%======================================================================
\section{Gap Analysis and Inputs for Derivation}\label{sec:gaps}
%%======================================================================

\subsection{Resolved Items (from v1)}

\begin{enumerate}[nosep]
  \item \textbf{Finding B} (incorrect projector-to-potential map):
    \textsc{resolved.}  Correct values
    $P^{(2)}_{0000} = 2/3$, $P^{(0\text{-}s)}_{0000} = 1/3$
    established and verified against three references.
  \item \textbf{Finding G} (placeholder in spin-2 argument):
    \textsc{resolved.}  All momentum arguments explicit:
    $z = k^2/\Lambda^2$.
  \item \textbf{Finding J} (preferred-frame asserted, not derived):
    \textsc{resolved.}  Derived from gauge invariance of spectral
    action (Finding~\ref{find:sct-conservative}).
\end{enumerate}

\subsection{Remaining Gaps}

\begin{center}
\begin{tabular}{clp{5cm}l}
\toprule
\# & Gap & Description & Resolution \\
\midrule
G1 & Full nonlocal & Exact $\Phi(r)$, $\Psi(r)$ beyond Yukawa &
  Numerical quadrature \\
G2 & $\beta_{\mathrm{PPN}}$ & Requires cubic vertex
  ($\mathcal{O}(h^3)$ from NT-4b) &
  Expand to cubic order \\
G3 & Analytic continuation & Euclidean $\to$ Lorentzian for
  dynamic sources &
  Deferred to MR-1 \\
G4 & Ghost pole prescription & $\Pi_{\mathrm{TT}}$ zero at
  $z_0 \approx 2.415$ &
  Fakeon (Anselmi) \\
G5 & $\xi$-dependent scan & Full parameter space exploration &
  Numerical scan \\
G6 & Future experiments & GAIA, BepiColombo, LISA pathfinder &
  Data release \\
\bottomrule
\end{tabular}
\end{center}

\subsection{Key Inputs for the PPN-1 Derivation}

The following formulas are the essential inputs for the PPN-1
derivation document:

\begin{enumerate}[nosep]
  \item \textbf{$K_\Phi$}: Eq.~\eqref{eq:K-Phi} ---
    $K_\Phi = 4/(3\Pi_{\mathrm{TT}}) - 1/(3\Pi_{\mathrm{s}})$.
  \item \textbf{$K_\Psi$}: Eq.~\eqref{eq:K-Psi} ---
    $K_\Psi = 2/(3\Pi_{\mathrm{TT}}) + 1/(3\Pi_{\mathrm{s}})$.
  \item \textbf{$\gamma_{\mathrm{eff}}(z)$}: Eq.~\eqref{eq:gamma-eff} ---
    ratio $K_\Psi/K_\Phi$.
  \item \textbf{Yukawa masses}: Eq.~\eqref{eq:sct-masses} ---
    $m_2 = \sqrt{60/13}\,\Lambda$,
    $m_0 = \Lambda/\sqrt{6(\xi-1/6)^2}$.
  \item \textbf{Stelle benchmark}: Eqs.~\eqref{eq:Phi-Stelle}--\eqref{eq:Psi-Stelle}.
  \item \textbf{Solar-system estimate}: $m_2 r_\oplus
    \approx 2 \times 10^{46}$ for $\Lambda \sim M_{\mathrm{Pl}}$,
    giving $|\gamma - 1| \sim e^{-10^{46}} \approx 0$.
  \item \textbf{Cassini bound}: Eq.~\eqref{eq:cassini} ---
    $|\gamma-1| < 2.3\times 10^{-5}$.
  \item \textbf{E\"ot-Wash bound}: $\Lambda > 2.4$~meV (approximate);
    $\Lambda > 2.565$~meV (precise, from LT-3d).
  \item \textbf{Preferred-frame}: all zero
    (Finding~\ref{find:sct-conservative}).
  \item \textbf{vDVZ}: SCT avoids the discontinuity
    (Finding~\ref{find:sct-vdvz}).
  \item \textbf{$\beta$}: $= 1 + \mathcal{O}(e^{-mr})$ at
    Newtonian order; cubic vertex needed for precision.
  \item \textbf{Cross-checks}: Stelle (\S\ref{sec:stelle}),
    IDG (\S\ref{sec:nonlocal}), Brans--Dicke (\S\ref{sec:BD}).
\end{enumerate}


%%======================================================================
% References
%%======================================================================
\begin{thebibliography}{30}

\bibitem{Will14}
C.~M.~Will,
``The confrontation between general relativity and experiment,''
\textit{Living Rev.\ Rel.} \textbf{17} (2014) 4,
\href{https://arxiv.org/abs/1403.7377}{arXiv:1403.7377 [gr-qc]}.

\bibitem{Ste77}
K.~S.~Stelle,
``Renormalization of higher-derivative quantum gravity,''
\textit{Phys.\ Rev.\ D} \textbf{16} (1977) 953.

\bibitem{Ste78}
K.~S.~Stelle,
``Classical gravity with higher derivatives,''
\textit{Gen.\ Rel.\ Grav.} \textbf{9} (1978) 353.

\bibitem{BMS06}
T.~Biswas, A.~Mazumdar, W.~Siegel,
``Bouncing universes in string-inspired gravity,''
\textit{JCAP} \textbf{0603} (2006) 009,
\href{https://arxiv.org/abs/hep-th/0508194}{hep-th/0508194}.

\bibitem{EKM16}
J.~Edholm, A.~S.~Koshelev, A.~Mazumdar,
``Behavior of the Newtonian potential for ghost-free gravity
and singularity-free gravity,''
\textit{Phys.\ Rev.\ D} \textbf{94} (2016) 104033,
\href{https://arxiv.org/abs/1604.01989}{arXiv:1604.01989 [gr-qc]}.

\bibitem{BGKM12}
T.~Biswas, E.~Gerwick, T.~Koivisto, A.~Mazumdar,
``Towards singularity and ghost free theories of gravity,''
\textit{Phys.\ Rev.\ Lett.} \textbf{108} (2012) 031101,
\href{https://arxiv.org/abs/1110.5249}{arXiv:1110.5249 [gr-qc]}.

\bibitem{BD61}
C.~Brans, R.~H.~Dicke,
``Mach's principle and a relativistic theory of gravitation,''
\textit{Phys.\ Rev.} \textbf{124} (1961) 925.

\bibitem{vDV70}
H.~van~Dam, M.~Veltman,
``Massive and massless Yang-Mills and gravitational fields,''
\textit{Nucl.\ Phys.\ B} \textbf{22} (1970) 397.

\bibitem{Zak70}
V.~I.~Zakharov,
``Linearized gravitation theory and the graviton mass,''
\textit{JETP Lett.} \textbf{12} (1970) 312.

\bibitem{Ber03}
B.~Bertotti, L.~Iess, P.~Tortora,
``A test of general relativity using radio links with the Cassini
spacecraft,''
\textit{Nature} \textbf{425} (2003) 374.

\bibitem{VFL14}
A.~K.~Verma, A.~Fienga, J.~Laskar, H.~Manche, M.~Gastineau,
``Use of MESSENGER radioscience data to improve planetary
ephemeris and to test general relativity,''
\textit{Astron.\ Astrophys.} \textbf{561} (2014) A115.

\bibitem{LLP09}
S.~B.~Lambert, C.~Le~Poncin-Lafitte,
``Determination of the relativistic parameter $\gamma$ using very
long baseline interferometry,''
\textit{Astron.\ Astrophys.} \textbf{499} (2009) 331.

\bibitem{WTB04}
J.~G.~Williams, S.~G.~Turyshev, D.~H.~Boggs,
``Progress in lunar laser ranging tests of relativistic gravity,''
\textit{Phys.\ Rev.\ Lett.} \textbf{93} (2004) 261101.

\bibitem{Tou22}
P.~Touboul et al.\ (MICROSCOPE Collaboration),
``MICROSCOPE Mission: Final Results of the Test of the Equivalence
Principle,''
\textit{Phys.\ Rev.\ Lett.} \textbf{129} (2022) 121102.

\bibitem{Lee20}
J.~G.~Lee, E.~G.~Adelberger, T.~S.~Cook, S.~M.~Fleischer,
B.~R.~Heckel,
``New test of the gravitational $1/r^2$ law at separations
down to $52\;\mu$m,''
\textit{Phys.\ Rev.\ Lett.} \textbf{124} (2020) 101101.

\bibitem{Kap07}
D.~J.~Kapner, T.~S.~Cook, E.~G.~Adelberger, J.~H.~Gundlach,
B.~R.~Heckel, C.~D.~Hoyle, H.~E.~Swanson,
``Tests of the gravitational inverse-square law below the
dark-energy length scale,''
\textit{Phys.\ Rev.\ Lett.} \textbf{98} (2007) 021101.

\bibitem{Tan20}
W.-H.~Tan et al.,
``Improvement for testing the gravitational inverse-square law
at the submillimeter range,''
\textit{Phys.\ Rev.\ Lett.} \textbf{124} (2020) 051301.

\bibitem{vN73}
P.~van~Nieuwenhuizen,
``On ghost-free tensor lagrangians and linearized gravitation,''
\textit{Nucl.\ Phys.\ B} \textbf{60} (1973) 478.

\bibitem{Riv64}
R.~J.~Rivers,
``Lagrangian theory for neutral massive spin-2 fields,''
\textit{Nuovo Cim.} \textbf{34} (1964) 386.

\bibitem{JP24}
D.~Jimu, T.~Prokopec,
``Gravitational potentials from quantum self-energy corrections,''
(2024),
\href{https://arxiv.org/abs/2410.01449}{arXiv:2410.01449 [gr-qc]}.

\bibitem{Acc02}
A.~Accioly, S.~Ragusa, H.~Mukai, E.~de Rey Neto,
``Algorithm for computing the propagator of higher-derivative
gravity theories,''
\textit{Int.\ J.\ Theor.\ Phys.} \textbf{39} (2000) 1599.

\bibitem{BV87}
A.~O.~Barvinsky, G.~A.~Vilkovisky,
``Beyond the Schwinger-DeWitt technique: Converting loops into trees
and in-in currents,''
\textit{Nucl.\ Phys.\ B} \textbf{282} (1987) 163.

\bibitem{Ans17}
D.~Anselmi,
``On the quantum field theory of the gravitational interactions,''
\textit{JHEP} \textbf{06} (2017) 086,
\href{https://arxiv.org/abs/1704.07728}{arXiv:1704.07728 [hep-th]}.

\bibitem{BCKM13}
T.~Biswas, A.~Conroy, A.~S.~Koshelev, A.~Mazumdar,
``Generalized ghost-free quadratic curvature gravity,''
\textit{Class.\ Quant.\ Grav.} \textbf{31} (2014) 015022,
\href{https://arxiv.org/abs/1308.2319}{arXiv:1308.2319 [hep-th]}.

\bibitem{MR15}
H.~Murata, M.~Tanaka,
``A review of short-range gravity experiments in the LHC era,''
\textit{Class.\ Quant.\ Grav.} \textbf{32} (2015) 033001,
\href{https://arxiv.org/abs/1408.3588}{arXiv:1408.3588 [hep-ex]}.

\bibitem{WTB12}
J.~G.~Williams, S.~G.~Turyshev, D.~H.~Boggs,
``Lunar laser ranging tests of the equivalence principle,''
\textit{Class.\ Quant.\ Grav.} \textbf{29} (2012) 184004.

\end{thebibliography}

\end{document}
